\documentclass[pract]{SCWorks}
\usepackage{preamble}
\setminted[py]{fontsize=\small, breaklines=true, style=bw, linenos}
\newcommand{\eqdef}{\stackrel {\rm def}{=}}
\newcommand{\specialcell}[2][c]{%
    \begin{tabular}[#1]{@{}c@{}}#2\end{tabular}}

\renewcommand\theFancyVerbLine{\small\arabic{FancyVerbLine}}

\newtheorem{lem}{Лемма}
\linespread{1.5}
\setminted{
    fontsize=\small,          % Размер шрифта
    baselinestretch=1.0,      % Единичный интервал
    breaklines=true,          % Перенос строк
    framesep=2mm,             % Отступ вокруг кода
    linenos=false,            % Нумерация строк (опционально)
    xleftmargin=2em,          % Отступ слева
}

\begin{document}
    \studentfemale


    \chair{информатики и программирования}

    \title{Методология разработки решений на основе искусственного интеллекта}

    \course{3}
    \term{5}

    \duration{14}
    \practStart{09.09}
    \practFinish{16.12}

    \group{341}
    \napravlenie{02.03.01 "--- Математическое обеспечение и администрирование информационных систем}
    \author{Хрустицкой Софьи Кирилловны}

    \chtitle{доцент, к.\,ф.-м.\,н.}
    \chname{М.\,В.\,Огнева}

    \satitle{доцент, к.\,пед.\,н.}
    \saname{А.\,В.\,Букушева}

    \date{2026}

    \newpage

    \maketitle
    \secNumbering
    \tableofcontents

    \section{Решение задач}
    В данном разделе представлены решения задач, выполненных в ходе учебной практики.
    Практические задания направлены на закрепление теоретического материала по линейной алгебре и отработку навыков применения её методов в вычислительной среде.
    
    \subsection{Задания 07.10.2025}
    \textbf{Условие задачи «Пропущенный KNN»}:

    Вася решает задачу классификации на два класса с помощью алгоритма KNN с количеством соседей 1.
    В обучающем датасете Васи 4 точки.
    Вася построил разделяющую поверхность, которая получилась у его алгоритма (см. рисунок \ref{fig:knn}).
    Красные точки плоскости алгоритм относит к классу 0, зеленые точки плоскости алгоритм относит к классу 1.
    Пока Вася пил чай, в компьютере произошел сбой, и один элементов датасета пропал.
    Помогите Васе восстановить четвертый элемент $x_4$, зная первые три элемента и координаты точек $a_1$ и $a_2$:
    \begin{itemize}
        \item[$\bullet$] $x_1 = (-2, 3)$
        \item[$\bullet$] $x_2 = (2, 8)$
        \item[$\bullet$] $x_3 = (1, 1)$
        \item[$\bullet$] $a_1 = (17/18, 25/6)$
        \item[$\bullet$] $a_2 = (15/22, 109/22)$
    \end{itemize}

    \begin{figure}[H]
        \centering
        \includegraphics[width=0.8\textwidth]{images/1.png}
        \caption{Разделяющая поверхность Васи}
        \label{fig:knn}
    \end{figure}

    \textbf{Решение}:

    Для алгоритма 1-NN граница между классами состоит из кусочков перпендикулярных серединных прямых между парами обучающих точек с разными классами.
    Точки, в которых сходятся сразу три такие границы, — это точки, равноудалённые от трёх обучающих объектов (вершины диаграммы Вороного).

    На рисунке \ref{fig:knn} голубые точки $a_1$ и $a_2$ лежат на границе между красной и зелёной областями и являются именно такими вершинами.
    Значит $a_1$ равноудалена от трёх точек $x_1$, $x_3$ и $x_4$, а $a_2$ равноудалена от трёх точек $x_1$, $x_2$ и $x_4$.
    Это и даёт уравнения для поиска $x_4$.

    Найдем расстояния от $a_1$ и $a_2$ до известных точек.
    Посчитаем квадраты расстояний от $a_1$ до $x_1$ и $x_3$:
    \[
    \begin{aligned}
    d^2(a_1,x_1) &= \left(\frac{17}{18} - (-2)\right)^2 + \left(\frac{25}{6} - 3\right)^2 \\
    &= \left(\frac{53}{18}\right)^2 + \left(\frac{7}{6}\right)^2 \\
    &= \frac{2809}{324} + \frac{49}{36} \\
    &= \frac{2809}{324} + \frac{441}{324} \\
    &= \frac{3250}{324} = \frac{1625}{162},
    \end{aligned}
    \]

    \[
    \begin{aligned}
    d^2(a_1,x_3) &= \left(\frac{17}{18} - 1\right)^2 + \left(\frac{25}{6} - 1\right)^2 \\
    &= \left(-\frac{1}{18}\right)^2 + \left(\frac{19}{6}\right)^2 \\
    &= \frac{1}{324} + \frac{361}{36} \\
    &= \frac{1}{324} + \frac{3249}{324} \\
    &= \frac{3250}{324} = \frac{1625}{162}.
    \end{aligned}
    \]

    Действительно,
    \[
    d^2(a_1,x_1) = d^2(a_1,x_3) = \frac{1625}{162}.
    \]

    Значит, так как $a_1$ — вершина диаграммы, расстояние от неё до $x_4$ такое же:

    \begin{equation}
    d^2(a_1,x_4) = \frac{1625}{162}.
    \end{equation}

    Теперь для точки $a_2$:

    \[
    \begin{aligned}
    d^2(a_2,x_1) &= \left(\frac{15}{22} - (-2)\right)^2 + \left(\frac{109}{22} - 3\right)^2 \\
    &= \left(\frac{59}{22}\right)^2 + \left(\frac{43}{22}\right)^2 \\
    &= \frac{3481}{484} + \frac{1849}{484} \\
    &= \frac{5330}{484} = \frac{2665}{242},
    \end{aligned}
    \]

    \[
    \begin{aligned}
    d^2(a_2,x_2) &= \left(\frac{15}{22} - 2\right)^2 + \left(\frac{109}{22} - 8\right)^2 \\
    &= \left(-\frac{29}{22}\right)^2 + \left(-\frac{67}{22}\right)^2 \\
    &= \frac{841}{484} + \frac{4489}{484} \\
    &= \frac{5330}{484} = \frac{2665}{242}.
    \end{aligned}
    \]

    Получаем
    \[
    d^2(a_2,x_1) = d^2(a_2,x_2) = \frac{2665}{242}.
    \]

    Следовательно, из того, что $a_2$ тоже вершина:
    \begin{equation}
    d^2(a_2,x_4) = \frac{2665}{242}.
    \end{equation}

    Составим системы для неизвестной точки $x_4=(u,v)$

    Из (1):
    \begin{equation}
    \left(u - \frac{17}{18}\right)^2 + \left(v - \frac{25}{6}\right)^2 = \frac{1625}{162}.
    \end{equation}

    Из (2):
    \begin{equation}
    \left(u - \frac{15}{22}\right)^2 + \left(v - \frac{109}{22}\right)^2 = \frac{2665}{242}.
    \end{equation}

    Это уравнения двух окружностей с центрами в $a_1$ и $a_2$ и известными радиусами.
    Ищем их точки пересечения.

    Раскрываем скобки и вычитаем (3) из (4):
    \[
    \begin{aligned}
    &\left(u^2 - \frac{15}{11}u + \frac{225}{484} + v^2 - \frac{109}{11}v + \frac{11881}{484}\right) \\
    - &\left(u^2 - \frac{17}{9}u + \frac{289}{324} + v^2 - \frac{25}{3}v + \frac{625}{36}\right) = \frac{2665}{242} - \frac{1625}{162}.
    \end{aligned}
    \]

    Упрощая, получаем линейное уравнение:
    \[
    \left(\frac{17}{9} - \frac{15}{11}\right)u + \left(\frac{25}{3} - \frac{109}{11}\right)v + \left(\frac{289}{324} + \frac{625}{36} - \frac{225}{484} - \frac{11881}{484}\right) = 0.
    \]

    После вычислений:
    \begin{equation}
    \frac{52}{99}u + \frac{68}{33}v - \frac{3136}{99} = 0.
    \end{equation}

    Упрощая (5):
    \begin{equation}
    4u + 51v - 241 = 0.
    \end{equation}

    Решая систему из уравнений (3) и (6), получаем две точки:
    \[
    (u,v) = (-2, 3) \quad \text{и} \quad (u,v) = (4, 5).
    \]

    Первая точка — это уже известный нам объект $x_1$, следовательно, второй вариант и есть искомый четвёртый элемент.
    \[
    \boxed{x_4 = (4, 5)}.
    \]

    По рисунку \ref{fig:knn} видно, что эта точка лежит в зелёной области, то есть относится к тому же классу, что и $x_2$ и $x_3$, тогда как $x_1$ — единственная точка класса 0.

    \textbf{Ответ}: $x_4 = (4, 5)$.

    \textbf{Условие задачи «Решающие деревья для Феди»}:

    Никита — электротехник, сын которого, Федя, изучает машинное обучение.
    Никита придумал конструктор, имитирующий работу решающего дерева для решения задачи бинарной классификации точек на плоскости, где точка описывается признаками ($x_1$, $x_2$).
    Каждой вершине решающего дерева (включая листы) соответствует один элемент конструктора.
    Конструктор позволяет строить только деревья, решающие правила в вершинах которого имеют вид [$x_i > C$] для произвольной константы $C$.
    Федя позвонил папе и сообщил, что придумал датасет для классификации на 2 класса из 100 различных точек, причём координаты $x_1$, $x_2$  пробегают все возможные значения от 1 до 10. При этом Федя забыл сказать папе, какие точки к какому классу принадлежат.
    Никите нужно взять на работе достаточное количество элементов конструктора, чтобы Федя смог построить решающее дерево вне зависимости от того, каким классам принадлежат точки.
    Какое количество элементов ему понадобится?
    Иными словами, найдите такое минимальное число $N$, что, вне зависимости от разбиения точек на 2 класса, найдётся решающее дерево, идеально решающее эту задачу классификации, состоящее не более чем из $N$ вершин.

    \textbf{Решение}:

    Сначала надо показать, что всегда хватит не более чем $N$ вершин.
    Внутренние вершины задают сравнение по одной координате: $[x_1 > C]$ — вертикальное разделение, $[x_2 > C]$ — горизонтальное разделение.

    Каждая вершина делит текущий прямоугольник на два подпрямоугольника.
    В итоге каждое листовое правило соответствует ось-ориентированному прямоугольнику на сетке $\{1,\dots,10\}^2$.

    Построим дерево, которое разбивает всю область $\{1,\dots,10\}^2$ на $100$ прямоугольников $1 \times 1$ (по одной ячейке на каждую точку).
    Например, можно последовательно делить по $x_1$: $[x_1 > 1.5]$, $[x_1 > 2.5]$, ..., $[x_1 > 9.5]$ или по $x_2$: $[x_2 > 1.5]$, $[x_2 > 2.5]$, ..., $[x_2 > 9.5]$.

    В каждом листе остаётся ровно одна точка.
    Для любой разметки достаточно поменять метки в листьях (0 или 1), не меняя структуру дерева.

    Для дерева с $L = 100$ листьями (по одному на каждую точку) в двоичном дереве выполняется соотношение:
    \[
    L = I + 1,
    \]
    где $I$ — число внутренних вершин.

    Тогда:
    \[
    I = L - 1 = 100 - 1 = 99.
    \]

    Общее число вершин:
    \[
    N_{\text{верх}} = L + I = 100 + 99 = 199.
    \]

    Таким образом, точно достаточно 199 элементов конструктора.

    Покажем, что существует разметка, для которой любое дерево должно иметь не менее 199 вершин.
    Рассмотрим разметку по чётности суммы координат:
    \[
    y(x_1,x_2) = (x_1 + x_2) \bmod 2.
    \]

    То есть если $x_1 + x_2$ чётно — класс 0 (белая клетка), если $x_1 + x_2$ нечётно — класс 1 (чёрная клетка).

    Это обычная «шахматная доска» на сетке $10 \times 10$.

    Любой ось-ориентированный прямоугольник, содержащий не менее двух точек, будет содержать точки обоих классов:
    \begin{itemize}
        \item[$\bullet$] Если ширина по $x_1$ хотя бы 2: есть точки $(x_1, x_2)$ и $(x_1+1, x_2)$. Их суммы координат отличаются на 1, значит классы разные.
        \item[$\bullet$] Если ширина 1, но высота по $x_2$ хотя бы 2: есть точки $(x_1, x_2)$ и $(x_1, x_2+1)$. Их суммы также отличаются на 1, классы разные.
    \end{itemize}

    Следовательно: любой прямоугольник с хотя бы двумя точками содержит точки обоих классов.

    Чтобы области листьев были однородными по классу нужно чтобы каждый лист мог содержать не более одной точки.

    Тогда число листьев $L \ge 100$, число внутренних вершин $I = L - 1 \ge 99$ и общее число вершин $N_{\text{низ}} = L + I \ge 100 + 99 = 199$.
    Таким образом, для шахматной разметки любое корректное дерево должно иметь не менее 199 вершин.

    Следовательно, минимальное количество элементов конструктора, которое Никита должен взять, равно:
    \[
    \boxed{N = 199}.
    \]

    \textbf{Ответ}: $N = 199$.

    \textbf{Условие задачи «Хитрый интеграл»}:

    Какое наибольшее значение может принимать функционал
    \[
    I(p) = \int_{0}^{\infty} \exp\left(\frac{(x-4)^2 - (x+4)^2}{4}\right) \ln p(x)\, dx
    \]
    при условии:
    \[
    p(x) > 0, \quad \int_{0}^{\infty} p(x)\, dx = 1?
    \]

    \textbf{Решение}:

    Упростим показатель экспоненты:
    \[
    \begin{aligned}
    (x-4)^2 - (x+4)^2 &= (x^2 - 8x + 16) - (x^2 + 8x + 16) \\
    &= -16x.
    \end{aligned}
    \]

    Тогда:
    \[
    \frac{(x-4)^2 - (x+4)^2}{4} = -4x.
    \]

    Следовательно:
    \[
    \exp\left(\frac{(x-4)^2 - (x+4)^2}{4}\right) = e^{-4x}.
    \]

    Таким образом, функционал принимает вид:
    \begin{equation}
    I(p) = \int_{0}^{\infty} e^{-4x} \ln p(x)\, dx.
    \end{equation}

    Максимизируем функционал (7) при условии нормировки:
    \[
    \int_{0}^{\infty} p(x)\, dx = 1.
    \]

    Составим расширенный функционал с множителем Лагранжа $\lambda$:
    \[
    J[p] = \int_{0}^{\infty} \left[e^{-4x} \ln p(x) + \lambda p(x)\right] dx - \lambda.
    \]

    Вариационная производная по $p(x)$:
    \[
    \frac{\delta J}{\delta p(x)} = \frac{e^{-4x}}{p(x)} + \lambda.
    \]

    В точке максимума:
    \begin{equation}
    \frac{e^{-4x}}{p(x)} + \lambda = 0.
    \end{equation}

    Из (8) получаем:
    \[
    p(x) = -\frac{e^{-4x}}{\lambda}.
    \]

    Так как $p(x) > 0$, то $\lambda < 0$. Обозначим $c = -\lambda > 0$:
    \begin{equation}
    p(x) = \frac{1}{c} e^{-4x}.
    \end{equation}

    Из условия нормировки:
    \[
    1 = \int_{0}^{\infty} p(x)\, dx = \frac{1}{c} \int_{0}^{\infty} e^{-4x}\, dx.
    \]

    Вычислим интеграл:
    \[
    \int_{0}^{\infty} e^{-4x}\, dx = \left[-\frac{1}{4} e^{-4x}\right]_{0}^{\infty} = \frac{1}{4}.
    \]

    Следовательно:
    \[
    1 = \frac{1}{c} \cdot \frac{1}{4} \quad \Rightarrow \quad c = \frac{1}{4}.
    \]

    Подставляя в (9), получаем оптимальную функцию:
    \begin{equation}
    \boxed{p^*(x) = 4e^{-4x}, \quad x \ge 0}.
    \end{equation}

    Это экспоненциальное распределение с параметром $\lambda = 4$.

    Подставим $p^*(x)$ в функционал (7):
    \[
    \begin{aligned}
    I_{\max} &= I(p^*) = \int_{0}^{\infty} e^{-4x} \ln\left(4e^{-4x}\right) dx \\
    &= \int_{0}^{\infty} e^{-4x} (\ln 4 - 4x)\, dx \\
    &= \ln 4 \int_{0}^{\infty} e^{-4x}\, dx - 4 \int_{0}^{\infty} x e^{-4x}\, dx.
    \end{aligned}
    \]

    Первый интеграл:
    \[
    \int_{0}^{\infty} e^{-4x}\, dx = \frac{1}{4}.
    \]

    Второй интеграл (интегрируем по частям или используем формулу для момента экспоненциального распределения):
    \[
    \int_{0}^{\infty} x e^{-4x}\, dx = \frac{1}{4^2} = \frac{1}{16}.
    \]

    \[
    \begin{aligned}
    I_{\max} &= \ln 4 \cdot \frac{1}{4} - 4 \cdot \frac{1}{16} \\
    &= \frac{\ln 4}{4} - \frac{1}{4} \\
    &= \frac{\ln 4 - 1}{4}.
    \end{aligned}
    \]

    В альтернативной форме (учитывая, что $\ln 4 = 2\ln 2$):

    \[
    I_{\max} = \frac{2\ln 2 - 1}{4}.
    \]

    \textbf{Ответ}:
    \[
    \boxed{I_{\max} = \frac{\ln 4 - 1}{4} = 0.097}.
    \]

    \textbf{Условие задачи «Неоднозначный PCA»}:

    Известно, что алгоритм PCA не всегда даёт однозначный ответ, потому что максимальный разброс может достигаться в проекции на несколько различных направляющих векторов.
    Дан набор данных из 8 точек в трёхмерном пространстве, где точки являются вершинами 3-мерного куба.
    Для этого набора данных обучили PCA с одной главной компонентой (n\_components=1 в библиотеке sklearn).
    Какое количество различных ответов могло получиться?
    Иными словами, найдите количество направлений, проекции точек 3-мерного куба на которые имеют максимальный разброс.
    Направления, отличающиеся только знаком направляющего вектора, считаются различными.

    \textbf{Решение}:

    Алгоритм PCA с n\_components = 1 выполняет следующие шаги:
    \begin{enumerate}
        \item[$\bullet$] Центрирование данных – вычитание среднего значения по каждому признаку.
        \item[$\bullet$] Вычисление ковариационной матрицы:
        \[
        \Sigma = \frac{1}{n} \sum_{k=1}^{n} (x^{(k)} - \bar{x})(x^{(k)} - \bar{x})^{\top}.
        \]
        \item[$\bullet$] Нахождение собственных значений и собственных векторов $\Sigma$.
        \item[$\bullet$] Выбор собственного вектора, соответствующего наибольшему собственному значению, в качестве первой главной компоненты.
    \end{enumerate}

    Пусть точки — вершины единичного куба $\{0,1\}^3$. Среднее значение каждой координаты:
    \[
    \bar{x}_1 = \bar{x}_2 = \bar{x}_3 = \frac{1}{2}.
    \]

    После центрирования каждая координата принимает значения $\pm\frac{1}{2}$ с равной вероятностью. Таким образом, имеем 8 точек вида:
    \[
    \left(\pm\frac{1}{2},\ \pm\frac{1}{2},\ \pm\frac{1}{2}\right),
    \]
    где все комбинации знаков равновероятны.

    Вычислим элементы ковариационной матрицы $\Sigma$.
    Каждая координата $X_i$ принимает значения $\pm\frac{1}{2}$ с вероятностью $\frac{1}{2}$:
    \[
    \operatorname{Var}(X_i) = E[X_i^2] = \left(\frac{1}{2}\right)^2 = \frac{1}{4}.
    \]

    Разные координаты независимы и симметричны относительно нуля:
    \[
    \operatorname{Cov}(X_i, X_j) = E[X_i X_j] = 0 \quad \text{при } i \neq j.
    \]

    Таким образом, ковариационная матрица имеет вид:
    \[
    \Sigma =
    \begin{pmatrix}
    \frac{1}{4} & 0 & 0 \\
    0 & \frac{1}{4} & 0 \\
    0 & 0 & \frac{1}{4}
    \end{pmatrix}
    = \frac{1}{4} I_3,
    \label{eq:covariance_matrix}
    \]
    где $I_3$ — единичная матрица размера $3\times 3$.

    Для матрицы $\Sigma = \frac{1}{4} I_3$:
    \begin{itemize}
        \item[$\bullet$] Все три собственных значения равны:
        \[
        \lambda_1 = \lambda_2 = \lambda_3 = \frac{1}{4}.
        \]
        \item[$\bullet$] Максимальное собственное значение имеет кратность 3.
        \item[$\bullet$] Любой единичный вектор $u \in \mathbb{R}^3$ ($\|u\| = 1$) является собственным вектором:
        \[
        \Sigma u = \frac{1}{4} I_3 u = \frac{1}{4} u.
        \]
    \end{itemize}

    В ситуации, когда максимальное собственное значение кратно, PCA с n\_components = 1 может вернуть любой единичный вектор.
    На практике конкретный выбор зависит от численных особенностей реализации алгоритма, cпособа решения задачи на собственные значения и начальных приближений в итерационных методах.

    Все направления дают одинаковый максимальный разброс проекций:
    \[
    \operatorname{Var}(u^{\top}X) = u^{\top}\Sigma u = \frac{1}{4} \|u\|^2 = \frac{1}{4}.
    \]

    Направление главной компоненты — это единичный вектор на сфере:
    \[
    S^2 = \{u \in \mathbb{R}^3 : \|u\| = 1\}.
    \]

    Согласно условию задачи, векторы $u$ и $-u$ считаются различными.

    Множество всех единичных векторов $S^2$ имеет мощность континуума (несчётное бесконечное множество).
    Учёт противоположных направлений как различных не изменяет мощность множества — оно остаётся континуумом.

    \textbf{Ответ}: Ковариационная матрица для вершин единичного куба пропорциональна единичной матрице, поэтому максимальное собственное значение имеет кратность 3, и любое направление является допустимой главной компонентой.

    \textbf{Условие задачи «Метод ближайших соседей»}:

    Рассмотрим $l$ точек, распределенных равномерно по объему $n$-мерного единичного шара с центром в нуле.
    Предположим, что мы хотим применить метод одного ближайшего соседа для точки $O$ — начала координат.
    Зададимся вопросом, на каком расстоянии будет расположен ближайший объект к точке $O$.
    Для ответа на этот вопрос выведите выражение для медианы расстояния от $O$ до ближайшего объекта.
    Дайте ответ для $l = 500$ и $n = 10$.

    \textbf{Решение}:

    В $n$-мерном единичном шаре равномерно распределены $\ell$ независимых точек.
    Требуется найти медиану расстояния от начала координат $O$ до ближайшей из этих точек.
    Рассмотреть частный случай при $n=10$, $\ell=500$.

    Пусть $X \in \mathbb{R}^n$ — случайная точка, равномерно распределённая в единичном $n$-мерном шаре.
    Обозначим расстояние до начала координат:
    \[
    R = \|X\|.
    \]

    Объём $n$-мерного шара радиуса $r$ пропорционален $r^n$:
    \[
    V_n(r) = c_n r^n,
    \]
    где $c_n$ — константа, зависящая от размерности.

    Функция распределения расстояния $R$:
    \[
    F_R(r) = \mathbb{P}(R \le r) = \frac{V_n(r)}{V_n(1)} = \frac{c_n r^n}{c_n \cdot 1^n} = r^n, \quad 0 \le r \le 1.
    \]

    Таким образом, плотность распределения:
    \[
    f_R(r) = \frac{d}{dr} F_R(r) = n r^{n-1}, \quad 0 \le r \le 1.
    \]

    Пусть $R_1, R_2, \dots, R_\ell$ — расстояния до начала координат для $\ell$ независимых точек.
    Рассмотрим минимум:
    \[
    R_{(1)} = \min_{i=1,\dots,\ell} R_i.
    \]

    Функция распределения минимума:
    \[
    \begin{aligned}
    F_{R_{(1)}}(r) &= \mathbb{P}(R_{(1)} \le r) = 1 - \mathbb{P}(R_{(1)} > r) \\
    &= 1 - \mathbb{P}(R_1 > r, \dots, R_\ell > r) \\
    &= 1 - \prod_{i=1}^{\ell} \mathbb{P}(R_i > r) \\
    &= 1 - \bigl(1 - F_R(r)\bigr)^{\ell} \\
    &= 1 - (1 - r^n)^{\ell}, \quad 0 \le r \le 1.
    \end{aligned}
    \]

    Плотность распределения минимума:
    \[
    f_{R_{(1)}}(r) = \frac{d}{dr} F_{R_{(1)}}(r) = \ell n r^{n-1} (1 - r^n)^{\ell-1}, \quad 0 \le r \le 1.
    \]

    Медиана $m$ определяется условием:
    \[
    F_{R_{(1)}}(m) = \frac{1}{2}.
    \]

    Подставляем выражение для функции распределения:
    \[
    1 - (1 - m^n)^{\ell} = \frac{1}{2}.
    \]

    Решаем уравнение:
    \[
    \begin{aligned}
    (1 - m^n)^{\ell} &= \frac{1}{2}, \\
    1 - m^n &= 2^{-1/\ell}, \\
    m^n &= 1 - 2^{-1/\ell}.
    \end{aligned}
    \]

    Отсюда получаем общую формулу для медианы:
    \[
    \boxed{m(\ell, n) = \bigl(1 - 2^{-1/\ell}\bigr)^{1/n}}.
    \]

    Подставляем значения в формулу:
    \[
    \begin{aligned}
    m(500, 10) &= \bigl(1 - 2^{-1/500}\bigr)^{1/10}.
    \end{aligned}
    \]

    Вычислим последовательно:
    1. $2^{-1/500} = e^{-\frac{\ln 2}{500}} \approx e^{-0.00138629} \approx 0.998614$.
    2. $1 - 2^{-1/500} \approx 1 - 0.998614 = 0.001386$.
    3. $(0.001386)^{1/10} = (0.001386)^{0.1}$.

    Или в более точном виде:
    \[
    \begin{aligned}
    2^{-1/500} &\approx 0.998614, \\
    1 - 2^{-1/500} &\approx 0.001386, \\
    m &\approx (0.001386)^{0.1} \approx 0.5178.
    \end{aligned}
    \]

    Для больших $\ell$ можно использовать приближение:
    \[
    2^{-1/\ell} = e^{-\frac{\ln 2}{\ell}} \approx 1 - \frac{\ln 2}{\ell}.
    \]

    Тогда:
    \[
    m^n \approx \frac{\ln 2}{\ell}, \quad m \approx \left(\frac{\ln 2}{\ell}\right)^{1/n}.
    \]

    Для $\ell=500$, $n=10$:
    \[
    m \approx \left(\frac{0.693147}{500}\right)^{0.1} = (0.00138629)^{0.1} \approx 0.5178,
    \]
    что совпадает с точным вычислением.

    \textbf{Ответ}: $m = 0.5$.

    \textbf{Условие задачи «Градиентный спуск»}:

    Пусть мы ищем минимум функции $f(x) = x^4$ с помощью градиентного спуска.
    Стартовая точка $x_0 = 4$. Используется обычный градиентный спуск, шаг которого выражается формулой
    \[
        x_{i+1} = x_i - \alpha \cdot \frac{df(x)}{dx}(x_i).
    \]
    При каком минимальном значении $\alpha > 0$ алгоритм не сойдется?
    Иными словами, при каком минимальном значении $\alpha > 0$ будет существовать такое $\delta > 0$, то алгоритм градиентного спуска, стартуя из точки $x_0$, никогда не придёт в интервал ($-\delta, \delta$)?

    \textbf{Решение}:

    Производная функции:
    \[
    f'(x) = 4x^3.
    \]

    Шаг градиентного спуска:
    \begin{equation}
    x_{k+1} = x_k - \alpha f'(x_k) = x_k - 4\alpha x_k^3 = x_k(1 - 4\alpha x_k^2).
    \end{equation}

    Рассмотрим модуль:
    \begin{equation}
    |x_{k+1}| = |x_k| \cdot |1 - 4\alpha x_k^2|.
    \end{equation}

    Пусть $0 < \alpha < \frac{1}{32}$ и $|x_k| \le 4$.
    Тогда:
    \[
    4\alpha x_k^2 \le 4\alpha \cdot 16 = 64\alpha < 64 \cdot \frac{1}{32} = 2.
    \]

    Следовательно:
    \[
    0 < 4\alpha x_k^2 < 2 \quad \Rightarrow \quad |1 - 4\alpha x_k^2| < 1.
    \]

    Из (12) получаем:
    \[
    |x_{k+1}| = |x_k| \cdot |1 - 4\alpha x_k^2| < |x_k|,
    \]
    при любом $x_k \neq 0$, $|x_k| \le 4$.

    \begin{itemize}
        \item[$\bullet$] При $x_0 = 4$: $|x_1| < 4$.
        \item[$\bullet$] По индукции: все $|x_k| \le 4$ и последовательность $|x_k|$ строго убывает.
        \item[$\bullet$] Последовательность ограничена снизу нулём, значит существует предел $L = \lim_{k \to \infty} x_k$.
    \end{itemize}

    Подставляя $L$ в (11):
    \[
    L = L - 4\alpha L^3 \quad \Rightarrow \quad 4\alpha L^3 = 0 \quad \Rightarrow \quad L = 0.
    \]

    Таким образом, при $0 < \alpha < \frac{1}{32}$ градиентный спуск сходится к нулю, и для любого $\delta > 0$ рано или поздно попадает в интервал $(-\delta, \delta)$.

    Рассмотрим $\alpha = \frac{1}{32}$:
    \[
    \begin{aligned}
    x_1 &= 4 - 4 \cdot \frac{1}{32} \cdot 4^3 = 4 - \frac{4 \cdot 64}{32} = 4 - 8 = -4, \\[4pt]
    x_2 &= -4 - 4 \cdot \frac{1}{32} \cdot (-4)^3 = -4 - 4 \cdot \frac{1}{32} \cdot (-64) = -4 + 8 = 4.
    \end{aligned}
    \]

    Получаем периодическую последовательность:
    \[
    x_0 = 4,\quad x_1 = -4,\quad x_2 = 4,\quad x_3 = -4,\quad \dots
    \]

    Период равен 2: $4 \leftrightarrow -4$.

    Последовательность никогда не заходит в интервал $(-4, 4)$. Можно взять $\delta = 4$.
    Следовательно, при $\alpha = \frac{1}{32}$ метод не сходится к минимуму.

    При $\alpha > \frac{1}{32}$:
    \[
    x_1 = 4 - 256\alpha < 4 - 8 = -4.
    \]

    Для $|x| \ge 4$ имеем:
    \[
    4\alpha x^2 \ge 4\alpha \cdot 16 > 2.
    \]

    Тогда:
    \[
    |1 - 4\alpha x^2| > 1 \quad \Rightarrow \quad |x_{k+1}| > |x_k|.
    \]

    Модуль начинает расти, и итерации удаляются от нуля. Такой алгоритм также не сходится.

    \textbf{Ответ}:
    \[
    \boxed{\alpha_{\min} = \dfrac{1}{32}}.
    \].

    \textbf{Условие задачи «Многоклассовая логистическая регрессия»}:

    Для задачи классификации на $k$ классов ${1, \dots, k}$ обучили алгоритм многоклассовой логистической регрессии, который устроен следующим образом.
    Если объекты в датасете имеют $n$ признаков, то обучаемыми параметрами алгоритма являются матрица $W$ размера $k \times n$ и вектор $b$ размера $k$.
    Вектор вероятностей принадлежности классам $1, \dots, k$ вычисляется по формуле
    \[
        p = softmax(Wx + b).
    \]
    На картинке \ref{fig:log_reg} изображено, как классификатор разбивает пространство признаков на три части, соответствующие классам.
    Здесь $n = 2$, $k = 3$, \[
    W = \begin{pmatrix}
    -3 & 3 \\
    1 & -2 \\
    2 & -1
    \end{pmatrix}, \quad
    b = \begin{pmatrix}
    8 \\ 1 \\ -9
    \end{pmatrix}.
    \]
    Пусть $A$ – точка пересечений всех трёх линий разделения классов. Найдите координаты точки $A$.

    \begin{figure}[H]
        \centering
        \includegraphics[width=0.8\textwidth]{images/2.png}
        \caption{Пространство признаков разбитое на три части}
        \label{fig:log_reg}
    \end{figure}

    \textbf{Решение}:

    Дана многоклассовая логистическая регрессия с тремя классами. Логиты вычисляются как:
    \[
    z = Wx + b,
    \]

    где
    \[
    W = \begin{pmatrix}
    -3 & 3 \\
    1 & -2 \\
    2 & -1
    \end{pmatrix}, \quad
    b = \begin{pmatrix}
    8 \\ 1 \\ -9
    \end{pmatrix}, \quad
    x = \begin{pmatrix}
    x_1 \\ x_2
    \end{pmatrix}.
    \]

    Обозначим строки матрицы $W$ как $w_1, w_2, w_3$, а компоненты вектора $b$ как $b_1, b_2, b_3$:
    \[
    \begin{aligned}
    w_1 &= (-3, 3), & b_1 &= 8, \\
    w_2 &= (1, -2), & b_2 &= 1, \\
    w_3 &= (2, -1), & b_3 &= -9.
    \end{aligned}
    \]

    Требуется найти координаты точки $A$, которая является пересечением всех трёх границ решений.

    Граница между классами $i$ и $j$ определяется уравнением:
    \[
    w_i^\top x + b_i = w_j^\top x + b_j.
    \]

    Точка $A$ принадлежит всем трём границам одновременно:
    \[
    w_1^\top x + b_1 = w_2^\top x + b_2 = w_3^\top x + b_3.
    \]

    Достаточно решить систему из двух независимых уравнений.
    \[
    w_1^\top x + b_1 = w_2^\top x + b_2.
    \]

    Подставляем:
    \[
    (-3, 3) \cdot \begin{pmatrix} x_1 \\ x_2 \end{pmatrix} + 8 = (1, -2) \cdot \begin{pmatrix} x_1 \\ x_2 \end{pmatrix} + 1.
    \]

    Раскрываем скалярные произведения:
    \[
    -3x_1 + 3x_2 + 8 = x_1 - 2x_2 + 1.
    \]

    Переносим все члены влево:
    \[
    \begin{aligned}
    -3x_1 + 3x_2 + 8 - x_1 + 2x_2 - 1 &= 0, \\
    -4x_1 + 5x_2 + 7 &= 0.
    \end{aligned}
    \]

    Получаем уравнение:
    \begin{equation}
    -4x_1 + 5x_2 = -7.
    \end{equation}

    \[
    w_1^\top x + b_1 = w_3^\top x + b_3.
    \]

    Подставляем:
    \[
    (-3, 3) \cdot \begin{pmatrix} x_1 \\ x_2 \end{pmatrix} + 8 = (2, -1) \cdot \begin{pmatrix} x_1 \\ x_2 \end{pmatrix} - 9.
    \]
    Раскрываем скалярные произведения:
    \[
    -3x_1 + 3x_2 + 8 = 2x_1 - x_2 - 9.
    \]

    Переносим все члены влево:
    \[
    \begin{aligned}
    -3x_1 + 3x_2 + 8 - 2x_1 + x_2 + 9 &= 0, \\
    -5x_1 + 4x_2 + 17 &= 0.
    \end{aligned}
    \]

    Получаем уравнение:
    \begin{equation}
    -5x_1 + 4x_2 = -17.
    \end{equation}

    Решаем систему из уравнений (13) и (14):
    \[
    \begin{cases}
    -4x_1 + 5x_2 = -7, \\
    -5x_1 + 4x_2 = -17.
    \end{cases}
    \]

    Из уравнения (13) выразим $x_2$:
    \begin{equation}
    5x_2 = 4x_1 - 7 \quad \Rightarrow \quad x_2 = \frac{4x_1 - 7}{5}.
    \end{equation}

    Подставляем (15) в уравнение (14):
    \[
    -5x_1 + 4 \cdot \frac{4x_1 - 7}{5} = -17.
    \]

    Умножаем обе части на 5:
    \[
    -25x_1 + 16x_1 - 28 = -85.
    \]

    Упрощаем:
    \[
    -9x_1 - 28 = -85,
    \]
    \[
    -9x_1 = -57,
    \]
    \[
    x_1 = \frac{57}{9} = \frac{19}{3}.
    \]

    Подставляем $x_1 = \frac{19}{3}$ в уравнение (15):
    \[
    \begin{aligned}
    x_2 &= \frac{4 \cdot \frac{19}{3} - 7}{5} \\
    &= \frac{\frac{76}{3} - 7}{5} \\
    &= \frac{\frac{76}{3} - \frac{21}{3}}{5} \\
    &= \frac{\frac{55}{3}}{5} \\
    &= \frac{55}{15} = \frac{11}{3}.
    \end{aligned}
    \]

    \textbf{Ответ}:
    \[
    \boxed{A = \left(\dfrac{19}{3},\; \dfrac{11}{3}\right)}.
    \]

    \textbf{Условие задачи «Мандаринки»}:

    Тася и ее друг Ваня очень любят мандарины, но Тася предпочитает абхазские, а Ваня — марокканские.
    Многочисленные эксперименты с мандаринами позволили ребятам вывести четыре основных показателя внешнего вида мандарина:
    \begin{itemize}
        \item[$\bullet$] толщина мандарина $x_1$ (численный признак);
        \item[$\bullet$] диаметр наибольшего сечения мандарина $x_2$ (численный признак);
        \item[$\bullet$] насыщенность цвета $x_3$ (численный признак);
        \item[$\bullet$] наличие/отсутствие листика $x_4$ (бинарный признак, принимает значение 0/1).
    \end{itemize}
    Численные признаки приведены в особых единицах, поэтому могут принимать как положительные, так и отрицательные значения.
    Более того, ребята установили, что распределение вектора параметров $(x_1, x_2, x_3, x_4)$ следующее: для обоих сортов вектор $(x_1, x_2, x_3)$ имеет трехмерное нормальное распределение, причём
     \begin{itemize}
        \item[$\bullet$] для абхазских манадринов — с вектором средних $a_1 = (0, 0 , 0)$ и матрицей ковариаций
        \[
        K_1 = \begin{pmatrix} 1 & 1 & 0 \\ 1 & 2 & 1 \\ 0 & 1 & 2 \end{pmatrix};
        \]
        \item[$\bullet$] для марокканских – с вектором средних $a_2 = (6, 6, 6)$ и матрицей ковариаций
        \[
        K_2 = \begin{pmatrix} 1 & 0 & 0 \\ 0 & 1 & 0 \\ 0 & 0 & 1 \end{pmatrix}.
        \]
    \end{itemize}
    Наконец, признак $x_4$ не зависит от совокупности признаков $(x_1, x_2, x_3)$ для обоих сортов, причём для абхазского имеет распределение $Bern(p_1)$, а для марокканского – $Bern(p_2)$ (здесь $p_i$ – вероятность наличия листика у мандаринки $i$-ого класса, $i = 1, 2$), $p_1 = 1/2$, $p_2 = 1/3$.
    Тася купила мандарины двух видов, оставила сумки на столе и пошла отдохнуть, но, пока она спала, кошка Мурка взобралась на стол, вывалила мандарины из обеих сумок и перемешала их.
    Помогите ребятам построить дискриминативную функцию методом максимального правдоподобия!

    \textbf{Решение}:

    Для класса абхазских мандаринов:
    \begin{itemize}
    \item[$\bullet$] Непрерывные признаки:
    \[
    \tilde{x} \mid y=1 \sim \mathcal{N}(a_1, K_1),
    \]
    где
    \[
    a_1 = \begin{pmatrix} 0 \\ 0 \\ 0 \end{pmatrix}, \quad
    K_1 = \begin{pmatrix}
        1 & 1 & 0 \\
        1 & 2 & 1 \\
        0 & 1 & 2
    \end{pmatrix}.
    \]

    \item[$\bullet$] Бинарный признак:
    \[
    x_4 \mid y=1 \sim \mathrm{Bern}(p_1), \quad p_1 = \frac{1}{2}.
    \]
    \end{itemize}

    Для класса марокканских мандаринов:

\begin{itemize}
    \item[$\bullet$] Непрерывные признаки:
    \[
    \tilde{x} \mid y=2 \sim \mathcal{N}(a_2, K_2),
    \]
    где
    \[
    a_2 = \begin{pmatrix} 6 \\ 6 \\ 6 \end{pmatrix}, \quad
    K_2 = \begin{pmatrix}
        1 & 0 & 0 \\
        0 & 1 & 0 \\
        0 & 0 & 1
    \end{pmatrix} = I_3.
    \]

    \item[$\bullet$] Бинарный признак:
    \[
    x_4 \mid y=2 \sim \mathrm{Bern}(p_2), \quad p_2 = \frac{1}{3}.
    \]
    \end{itemize}

    Априорные вероятности классов считаем равными: $P(y=1) = P(y=2) = \frac{1}{2}$.

    Плотность трёхмерного нормального распределения:
    \[
    \phi_i(\tilde{x}) = \frac{1}{(2\pi)^{3/2} \sqrt{\det K_i}}
    \exp\left(-\frac{1}{2} (\tilde{x} - a_i)^T K_i^{-1} (\tilde{x} - a_i)\right),
    \quad i = 1, 2.
    \]

    Плотность бинарного признака:
    \[
    \beta_i(x_4) = p_i^{x_4} (1 - p_i)^{1 - x_4}, \quad x_4 \in \{0, 1\}.
    \]

    В силу независимости признаков:
    \[
    p(x \mid y=i) = \phi_i(\tilde{x}) \cdot \beta_i(x_4), \quad i = 1, 2.
    \]

    В методе максимального правдоподобия (ММП) классификатор имеет вид:
    \[
    \hat{y} = \arg\max_{i=1,2} p(x \mid y=i).
    \]

    Введём дискриминативную функцию как логарифм отношения правдоподобий:
    \[
    g(x) = \ln\frac{p(x \mid y=1)}{p(x \mid y=2)}.
    \]

    Тогда решающее правило:
    \[
    \hat{y} =
    \begin{cases}
    1, & \text{если } g(x) > 0, \\
    2, & \text{если } g(x) < 0.
    \end{cases}
    \]

    Функция $g(x)$ распадается на сумму двух слагаемых:
    \[
    g(x) = g_{\text{Gauss}}(\tilde{x}) + g_{\text{leaf}}(x_4),
    \]
    где
    \[
    g_{\text{Gauss}}(\tilde{x}) = \ln\frac{\phi_1(\tilde{x})}{\phi_2(\tilde{x})}, \quad
    g_{\text{leaf}}(x_4) = \ln\frac{\beta_1(x_4)}{\beta_2(x_4)}.
    \]

    Вычислим $g_{\text{Gauss}}(\tilde{x})$:
    \[
    \begin{aligned}
    g_{\text{Gauss}}(\tilde{x}) &=
    -\frac{1}{2} \left[ (\tilde{x} - a_1)^T K_1^{-1} (\tilde{x} - a_1)
    - (\tilde{x} - a_2)^T K_2^{-1} (\tilde{x} - a_2) \right] \\
    &\quad + \frac{1}{2} \ln\frac{\det K_2}{\det K_1}.
    \end{aligned}
    \]

    Поскольку $\det K_1 = \det K_2 = 1$, логарифмический член обращается в ноль.

    Матрица $K_1^{-1}$:
    \[
    K_1^{-1} =
    \begin{pmatrix}
    3 & -2 & 1 \\
    -2 & 2 & -1 \\
    1 & -1 & 1
    \end{pmatrix}.
    \]

    Вычисляем квадратичные формы:
    \[
    \begin{aligned}
    (\tilde{x} - a_1)^T K_1^{-1} (\tilde{x} - a_1) &=
    \tilde{x}^T K_1^{-1} \tilde{x} \\
    &= 3x_1^2 - 4x_1x_2 + 2x_1x_3 + 2x_2^2 - 2x_2x_3 + x_3^2.
    \end{aligned}
    \]

    \[
    \begin{aligned}
    (\tilde{x} - a_2)^T K_2^{-1} (\tilde{x} - a_2) &=
    (\tilde{x} - a_2)^T (\tilde{x} - a_2) \\
    &= (x_1 - 6)^2 + (x_2 - 6)^2 + (x_3 - 6)^2 \\
    &= x_1^2 + x_2^2 + x_3^2 - 12(x_1 + x_2 + x_3) + 108.
    \end{aligned}
    \]

    Разность квадратичных форм:
    \[
    \begin{aligned}
    \tilde{x}^T &K_1^{-1} \tilde{x} - (\tilde{x} - a_2)^T (\tilde{x} - a_2) \\
    &= (3x_1^2 - 4x_1x_2 + 2x_1x_3 + 2x_2^2 - 2x_2x_3 + x_3^2) \\
    &\quad - (x_1^2 + x_2^2 + x_3^2 - 12x_1 - 12x_2 - 12x_3 + 108) \\
    &= 2x_1^2 - 4x_1x_2 + 2x_1x_3 + 12x_1 + x_2^2 - 2x_2x_3 + 12x_2 + 12x_3 - 108.
    \end{aligned}
    \]

    Таким образом:
    \begin{equation}
    g_{\text{Gauss}}(\tilde{x}) =
    -\frac{1}{2} \left( 2x_1^2 - 4x_1x_2 + 2x_1x_3 + 12x_1 + x_2^2 - 2x_2x_3 + 12x_2 + 12x_3 - 108 \right).
    \end{equation}

    Вычислим $g_{\text{leaf}}(x_4)$:
    \[
    \begin{aligned}
    g_{\text{leaf}}(x_4) &= \ln\frac{\beta_1(x_4)}{\beta_2(x_4)} \\
    &= \ln\left( \frac{p_1^{x_4} (1-p_1)^{1-x_4}}{p_2^{x_4} (1-p_2)^{1-x_4}} \right) \\
    &= x_4 \ln\frac{p_1}{p_2} + (1 - x_4) \ln\frac{1-p_1}{1-p_2}.
    \end{aligned}
    \]

    Подставляем $p_1 = \frac{1}{2}$, $p_2 = \frac{1}{3}$:
    \[
    \ln\frac{p_1}{p_2} = \ln\frac{1/2}{1/3} = \ln\frac{3}{2},
    \]
    \[
    \ln\frac{1-p_1}{1-p_2} = \ln\frac{1/2}{2/3} = \ln\frac{3}{4}.
    \]

    Тогда:
    \begin{equation}
    g_{\text{leaf}}(x_4) = x_4 \ln\frac{3}{2} + (1 - x_4) \ln\frac{3}{4}
    = \ln\frac{3}{4} + x_4 \ln 2.
    \end{equation}

    Объединяя (16) и (17), получаем:
    \begin{equation}
    \boxed{
    \begin{aligned}
    g(x_1, x_2, x_3, x_4) &=
    -\frac{1}{2} \Bigl( 2x_1^2 - 4x_1x_2 + 2x_1x_3 + 12x_1 \\
    &\quad + x_2^2 - 2x_2x_3 + 12x_2 + 12x_3 - 108 \Bigr) \\
    &\quad + \ln\frac{3}{4} + x_4 \ln 2.
    \end{aligned}
    }
    \end{equation}

    Классификация производится по знаку дискриминативной функции:
    \[
    \boxed{
    \hat{y} =
    \begin{cases}
    1 \text{ (абхазские)}, & \text{если } g(x_1, x_2, x_3, x_4) > 0, \\
    2 \text{ (марокканские)}, & \text{если } g(x_1, x_2, x_3, x_4) < 0.
    \end{cases}
    }
    \]

    Можно раскрыть скобки и привести подобные слагаемые:
    \[
    \begin{aligned}
    g(x) &= -x_1^2 + 2x_1x_2 - x_1x_3 - 6x_1 - \frac{1}{2}x_2^2 + x_2x_3 - 6x_2 \\
    &\quad - 6x_3 + 54 + \ln\frac{3}{4} + x_4 \ln 2.
    \end{aligned}
    \]

    Константы можно объединить:
    \[
    C = 54 + \ln\frac{3}{4} \approx 54 - 0.2877 = 53.7123.
    \]

    \textbf{Ответ}:
    \[
    \begin{aligned}
    g(x) &= -x_1^2 + 2x_1x_2 - x_1x_3 - 6x_1 - \frac{1}{2}x_2^2 + x_2x_3 - 6x_2 \\
    &\quad - 6x_3 + 54 + \ln\frac{3}{4} + x_4 \ln 2.
    \end{aligned}
    \]

    \subsection{Задания 21.10.2025}

    \textbf{Условие упражнения 4.1}:

    Напишите функцию Python, которая на входе принимает два вектора и на выходе выдает два числа: коэффициент корреляции Пирсона и значение косинусного сходства.
    Напишите исходный код, который следует формулам, представленным в данной главе; не используйте вызовы встроенной в NumPy функции np.corrcoef и встроенной в SciPy функции spatial.distance.cosine.
    Убедитесь, что два значения на выходе идентичны, когда переменные уже центрированы по среднему значению, и различны, когда переменные не центрированы по среднему значению.

    \textbf{Решение}:

    Для двух одномерных переменных $x = (x_1, \dots, x_n)$ и $y = (y_1, \dots, y_n)$ коэффициент корреляции Пирсона определяется как:
    \[
    \rho_{xy} = \frac{\sum_{i=1}^{n} (x_i - \bar{x})(y_i - \bar{y})}
    {\sqrt{\sum_{i=1}^{n} (x_i - \bar{x})^2} \cdot \sqrt{\sum_{i=1}^{n} (y_i - \bar{y})^2}},
    \]

    где:
    \begin{itemize}
        \item[$\bullet$] $\bar{x} = \frac{1}{n} \sum_{i=1}^{n} x_i$ — среднее значение $x$
        \item[$\bullet$] $\bar{y} = \frac{1}{n} \sum_{i=1}^{n} y_i$ — среднее значение $y$
    \end{itemize}

    В векторной форме, обозначив центрированные векторы:
    \[
    \tilde{x} = x - \bar{x}\mathbf{1}, \quad \tilde{y} = y - \bar{y}\mathbf{1},
    \]
    где $\mathbf{1}$ — вектор из единиц, формулу можно записать как:
    \[
    \rho_{xy} = \frac{\tilde{x} \cdot \tilde{y}}{\|\tilde{x}\| \cdot \|\tilde{y}\|}.
    \]

    Для двух векторов $x$ и $y$ косинусное сходство определяется как косинус угла между ними:
    \[
    \cos\theta = \frac{x \cdot y}{\|x\| \cdot \|y\|}.
    \]

    Согласно Коэну, если переменные уже центрированы, то корреляция Пирсона сводится к косинусному сходству.

    Реализуем функции корреляции Пирсона и косинусного сходства на Python. Будем использовать только библиотеку numpy.
    Первая функция принимает два вектора и возвращает кортеж, где на первом место стоит коэффициент корреляции Пирсона, а на втором косинусное сходство.

    \begin{minted}{python}
    def corr_and_cos(x, y):
        x = np.asarray(x, dtype=float)
        y = np.asarray(y, dtype=float)

        if x.shape != y.shape:
            raise ValueError("x и y должны иметь одинаковую длину")

        # --- Вычисление косинусного сходства ---
        dot_xy = np.sum(x * y)
        norm_x = np.sqrt(np.sum(x ** 2))
        norm_y = np.sqrt(np.sum(y ** 2))

        if norm_x == 0 or norm_y == 0:
            cosine = 0.0
        else:
            cosine = dot_xy / (norm_x * norm_y)

        # --- Вычисление корреляции Пирсона ---
        # Центрирование векторов
        x_c = x - x.mean()
        y_c = y - y.mean()

        # Скалярное произведение центрированных векторов
        dot_c = np.sum(x_c * y_c)

        # Нормы центрированных векторов
        norm_xc = np.sqrt(np.sum(x_c ** 2))
        norm_yc = np.sqrt(np.sum(y_c ** 2))

        # Избегаем деления на ноль
        if norm_xc == 0 or norm_yc == 0:
            pearson = 0.0
        else:
            pearson = dot_c / (norm_xc * norm_yc)

        return pearson, cosine
    \end{minted}

    Для удобства использования можно также реализовать отдельные функции для вычисления только одного значения из функции:
    \begin{minted}{python}
    def pearson_correlation(x, y):
        """Вычисляет только корреляцию Пирсона."""
        return corr_and_cos(x, y)[0]

    def cosine_similarity(x, y):
        """Вычисляет только косинусное сходство."""
        return corr_and_cos(x, y)[1]
    \end{minted}

    Сравним с библиотечными функциями. Используем библиотеки scipy и sklearn:
    \begin{minted}{python}
    def test_implementation():
        # Тестовые данные
        x = np.array([1, 2, 3, 4, 5])
        y = np.array([2, 4, 6, 8, 10])

        # Наша реализация
        pearson_our, cosine_our = corr_and_cos(x, y)

        # Библиотечная реализация корреляции
        pearson_lib, _ = stats.pearsonr(x, y)

        # Библиотечная реализация косинусного сходства
        cosine_lib = sk_cosine([x], [y])[0, 0]

        # Проверка
        print(f"Наша реализация: pearson = {pearson_our:.6f}, cosine = {cosine_our:.6f}")
        print(f"Библиотечная:   pearson = {pearson_lib:.6f}, cosine = {cosine_lib:.6f}")
        print(f"Совпадение pearson: {np.isclose(pearson_our, pearson_lib)}")
        print(f"Совпадение cosine:  {np.isclose(cosine_our, cosine_lib)}")

        return pearson_our, cosine_our
    \end{minted}
    Тест прошел успешно:
    \begin{verbatim}
        Наша реализация: pearson = 1.000000, cosine = 1.000000
        Библиотечная:   pearson = 1.000000, cosine = 1.000000
        Совпадение pearson: True
        Совпадение cosine:  True
    \end{verbatim}

    Проверим утверждения Коэна для центрированных данных, где корреляция Пирсона равна косиносному сходству.
    \begin{minted}{python}
    def test_cohen_assertion():
        # Исходные данные (нецентрированные)
        x = np.array([0, 1, 2, 3], dtype=float)
        y = np.array([100, 101, 102, 103], dtype=float)

        print("Тест 1: Нецентрированные данные")
        p1, c1 = corr_and_cos(x, y)
        print(f"  Pearson: {p1:.6f}")
        print(f"  Cosine:  {c1:.6f}")
        print(f"  Равны?: {np.isclose(p1, c1)}")

        # Центрированные данные
        xc = x - x.mean()
        yc = y - y.mean()

        print("\nТест 2: Центрированные данные")
        p2, c2 = corr_and_cos(xc, yc)
        print(f"  Pearson: {p2:.6f}")
        print(f"  Cosine:  {c2:.6f}")
        print(f"  Равны?: {np.isclose(p2, c2, atol=1e-10)}")

        return p1, c1, p2, c2
    \end{minted}
    Тесты прошли успешно:
    \begin{verbatim}
        Тест 1: Нецентрированные данные
        Pearson: 1.000000
        Cosine:  0.808317
        Равны?: False

        Тест 2: Центрированные данные
        Pearson: 1.000000
        Cosine:  1.000000
        Равны?: True
    \end{verbatim}

    \textbf{Условие упражнения 4.2}:

    Давайте продолжим обследовать разницу между корреляцией и косинусным сходством.
    Создайте переменную, содержащую целые числа от 0 до 3, вторую переменную, равную первой переменной плюс некоторое смещение.
    Затем создайте симуляцию, в которой вы систематически варьируете то смещение между –50 и +50 (то есть на первой итерации симуляции вторая переменная будет равна [–50, –49, –48, –47]).
    В цикле for вычислите корреляцию и косинусное сходство между двумя переменными и сохраните эти результаты.
    Затем постройте линейный график, показывающий, как среднее смещение влияет на корреляцию и косинусное сходство.
    Вы должны суметь воспроизвести рис. \ref{fig:res_4_2}.

    \begin{figure}[H]
        \centering
        \includegraphics[width=0.8\textwidth]{images/3.jpg}
        \caption{Результаты упражнения 4.2}
        \label{fig:res_4_2}
    \end{figure}

    \textbf{Решение}:

    Функция вычисления корреляции Пирсона и косинусного сходства была реализована в предыдущем задании.
    Сгенерируем данные и вычислим метрики:
    \begin{minted}{python}
    # Определение исходной переменной
    x = np.array([0, 1, 2, 3], dtype=float)

    # Определение диапазона сдвигов
    shifts = np.arange(-50, 51)  # [-50, -49, ..., 49, 50]

    # Инициализация массивов для результатов
    pearsons = []  # Для корреляций Пирсона
    cosines = []   # Для косинусных сходств

    # Цикл по всем значениям сдвига
    for d in shifts:
        y = x + d

        # Вычисление метрик
        r, c = corr_and_cos(x, y)

        # Сохранение результатов
        pearsons.append(r)
        cosines.append(c)

    # Преобразование в numpy массивы
    pearsons = np.array(pearsons)
    cosines = np.array(cosines)
    \end{minted}

    Визуализируем результаты с помощью библиотеки matplotplib:
    \begin{minted}{python}
    # Настройка параметров графика
    plt.figure(figsize=(10, 6))
    plt.rcParams['font.size'] = 12
    plt.rcParams['axes.linewidth'] = 1.5

    # Построение графика корреляции Пирсона
    plt.plot(shifts, pearsons,
             color='black',
             marker='o',
             markersize=4,
             linewidth=2,
             label='Корреляция Пирсона ($r$)')

    # Построение графика косинусного сходства
    plt.plot(shifts, cosines,
             color='gray',
             marker='s',
             markersize=4,
             linewidth=2,
             label='Косинусное сходство ($\cos\\theta$)')

    # Настройка осей и подписей
    plt.xlabel('Сдвиг $d$', fontsize=14)
    plt.ylabel('Значение метрики', fontsize=14)
    plt.title('Зависимость корреляции Пирсона и косинусного сходства от сдвига',
              fontsize=16, pad=20)

    # Добавление дополнительных элементов
    plt.axhline(y=1.0, color='red', linestyle='--', linewidth=1, alpha=0.5)
    plt.axhline(y=0.0, color='black', linestyle=':', linewidth=1, alpha=0.3)
    plt.axvline(x=0, color='black', linestyle=':', linewidth=1, alpha=0.3)

    # Легенда
    plt.legend(loc='upper right', fontsize=12, frameon=True)

    # Сетка
    plt.grid(True, linestyle='--', alpha=0.3)

    # Пределы по осям
    plt.xlim([-55, 55])
    plt.ylim([-1.1, 1.1])

    # Отображение графика
    plt.tight_layout()
    plt.show()
    \end{minted}

    Получился вот такой график (см. рис. \ref{fig:vis_4_2}), он идентичен графику из условия:

    \begin{figure}[H]
        \centering
        \includegraphics[width=0.8\textwidth]{images/4.jpg}
        \caption{График построенный для коэфф. Пирсона и косинуносного сход.}
        \label{fig:vis_4_2}
    \end{figure}

    Из рисунка \ref{fig:vis_4_2} можно понять, что корреляция Пирсона остаётся постоянной и равной 1 для всех значений сдвига.
    Косинусное сходство при $d = 0$: $\cos\theta = 1$ – векторы сонаправлены, при $d > 0$: $\cos\theta$ монотонно убывает, приближаясь к 0.2896 при $d = 50$ и при $d < 0$: $\cos\theta$ сначала уменьшается, достигает отрицательных значений, затем возвращается к 0.
    Только при $d = 0$ обе метрики равны.
    Корреляция Пирсона измеряет линейную зависимость между переменными.
    Поскольку $y = x + d$ — линейная функция, корреляция всегда равна 1.
    Косинусное сходство измеряет угол между векторами.
    При добавлении константы $d$ направление вектора меняется, поэтому косинус угла изменяется.
    Корреляция Пирсона инвариантна относительно сдвига (добавления константы).
    Косинусное сходство не инвариантно относительно сдвига.

    Для любых $d$ корреляция Пирсона равна 1:
    \[
    \rho(x, x + d) = 1
    \]

    Доказательство:
    \[
    \begin{aligned}
    y &= x + d \\
    \bar{y} &= \bar{x} + d \\
    y - \bar{y} &= (x + d) - (\bar{x} + d) = x - \bar{x} \\
    \rho(x, y) &= \frac{\sum (x_i - \bar{x})(y_i - \bar{y})}
    {\sqrt{\sum (x_i - \bar{x})^2} \sqrt{\sum (y_i - \bar{y})^2}} \\
    &= \frac{\sum (x_i - \bar{x})(x_i - \bar{x})}
    {\sqrt{\sum (x_i - \bar{x})^2} \sqrt{\sum (x_i - \bar{x})^2}} = 1
    \end{aligned}
    \]

    Косинусное сходство зависит от $d$:
    \[
    \cos(x, x + d) = \frac{\sum x_i(x_i + d)}
    {\sqrt{\sum x_i^2} \sqrt{\sum (x_i + d)^2}}
    \]

    Для нашего конкретного $x = [0, 1, 2, 3]$:
    \[
    \begin{aligned}
    \sum x_i &= 6, \quad \sum x_i^2 = 14 \\
    \cos(x, x + d) &= \frac{14 + 6d}{\sqrt{14} \sqrt{14 + 12d + 4d^2}}
    \end{aligned}
    \]

    Таблица ключевых значений:
    \begin{table}[H]
    \centering
    \begin{tabular}{|c|c|c|}
    \hline
    \textbf{Сдвиг $d$} & \textbf{Корреляция Пирсона} & \textbf{Косинусное сходство} \\
    \hline
    -50 & 1.0000 & -0.7526 \\
    -25 & 1.0000 & -0.1176 \\
    -10 & 1.0000 & 0.4856 \\
    0 & 1.0000 & 1.0000 \\
    10 & 1.0000 & 0.8083 \\
    25 & 1.0000 & 0.5144 \\
    50 & 1.0000 & 0.2896 \\
    \hline
    \end{tabular}
    \caption{Значения метрик для различных сдвигов}
    \label{tab:metrics}
    \end{table}

    \subsection{Задания 11.11.2025}

    \textbf{Условие задания 5.5}:

    Для трёх функций:
    \begin{enumerate}
        \item $f_1(x) = \sin x_1 \cos x_2$, где $x \in \mathbb{R}^2$
        \item $f_2(x, y) = x^\top y$, где $x, y \in \mathbb{R}^n$
        \item $f_3(x) = xx^\top$, где $x \in \mathbb{R}^n$
    \end{enumerate}

    Требуется:
    \begin{enumerate}[label=\alph*)]
        \item Указать размерности $\dfrac{\partial f_i}{\partial x}$
        \item Выписать явные выражения для якобианов
    \end{enumerate}

    \textbf{Решение}:

    Для функции $f: \mathbb{R}^m \to \mathbb{R}^k$ якобианом называется матрица:
    \[
    J_f(x) = \frac{\partial f}{\partial x} =
    \begin{bmatrix}
    \frac{\partial f_1}{\partial x_1} & \cdots & \frac{\partial f_1}{\partial x_m} \\
    \vdots & \ddots & \vdots \\
    \frac{\partial f_k}{\partial x_1} & \cdots & \frac{\partial f_k}{\partial x_m}
    \end{bmatrix} \in \mathbb{R}^{k \times m},
    \]
    где $f_i$ — $i$-я компонента функции $f$.

    Для матричной функции $F: \mathbb{R}^m \to \mathbb{R}^{p \times q}$ используется операция векторизации $\mathrm{vec}(\cdot)$, которая преобразует матрицу в вектор, выписывая её столбцы один за другим:
    \[
    \mathrm{vec}: \mathbb{R}^{p \times q} \to \mathbb{R}^{pq}.
    \]

    Решение для функции $f_1$:
    Функция $f_1: \mathbb{R}^2 \to \mathbb{R}$ имеет:
    \begin{itemize}
        \item[$\bullet$] Входная размерность: $m = 2$
        \item[$\bullet$] Выходная размерность: $k = 1$
    \end{itemize}

    Следовательно:
    \[
    \boxed{\frac{\partial f_1}{\partial x} \in \mathbb{R}^{1 \times 2}}.
    \]

    Вычислим частные производные:
    \[
    \begin{aligned}
    \frac{\partial f_1}{\partial x_1} &= \frac{\partial}{\partial x_1} (\sin x_1 \cos x_2) = \cos x_1 \cos x_2, \\
    \frac{\partial f_1}{\partial x_2} &= \frac{\partial}{\partial x_2} (\sin x_1 \cos x_2) = -\sin x_1 \sin x_2.
    \end{aligned}
    \]

    Якобиан как матрица-строка:
    \[
    \boxed{J_{f_1}(x) = \frac{\partial f_1}{\partial x} =
    \begin{bmatrix}
    \cos x_1 \cos x_2 & -\sin x_1 \sin x_2
    \end{bmatrix}}.
    \]

    Решение для функции $f_2$:
    Рассматриваем производную по $x$ при фиксированном $y$.

    Функция $f_2: \mathbb{R}^n \times \mathbb{R}^n \to \mathbb{R}$ по первому аргументу $x$:
    \begin{itemize}
        \item[$\bullet$] Входная размерность (по $x$): $m = n$
        \item[$\bullet$] Выходная размерность: $k = 1$
    \end{itemize}

    Следовательно:
    \[
    \boxed{\frac{\partial f_2}{\partial x} \in \mathbb{R}^{1 \times n}}.
    \]

    Распишем скалярное произведение:
    \[
    f_2(x, y) = \sum_{i=1}^n x_i y_i.
    \]

    Частные производные:
    \[
    \frac{\partial f_2}{\partial x_j} = \frac{\partial}{\partial x_j} \left( \sum_{i=1}^n x_i y_i \right) = y_j, \quad j = 1, \dots, n.
    \]

    Якобиан как матрица-строка:
    \[
    \boxed{J_{f_2}(x, y) = \frac{\partial f_2}{\partial x} =
    \begin{bmatrix}
    y_1 & y_2 & \cdots & y_n
    \end{bmatrix} = y^\top}.
    \]

    Дополнительно: производная по $y$.
    Аналогично:
    \[
    \boxed{\frac{\partial f_2}{\partial y} \in \mathbb{R}^{1 \times n}, \quad J_{f_2}^y(x, y) = \frac{\partial f_2}{\partial y} = x^\top}.
    \]

    Решение для функции $f_3$:
    Функция $f_3: \mathbb{R}^n \to \mathbb{R}^{n \times n}$ имеет матричный выход. Используем векторизацию:
    \[
    \tilde{f}_3(x) = \mathrm{vec}(xx^\top): \mathbb{R}^n \to \mathbb{R}^{n^2}.
    \]

    Тогда:
    \begin{itemize}
        \item[$\bullet$] Входная размерность: $m = n$
        \item[$\bullet$] Выходная размерность (после векторизации): $k = n^2$
    \end{itemize}

    Следовательно:
    \[
    \boxed{\frac{\partial f_3}{\partial x} \equiv \frac{\partial \mathrm{vec}(xx^\top)}{\partial x} \in \mathbb{R}^{n^2 \times n}}.
    \]

    Матрица $xx^\top$ имеет элементы:
    \[
    [xx^\top]_{ij} = x_i x_j, \quad i, j = 1, \dots, n.
    \]

    Рассмотрим функцию $\tilde{f}_{3,ij}(x) = x_i x_j$ как элемент векторизованной матрицы.

    Частные производные:
    \[
    \frac{\partial \tilde{f}_{3,ij}}{\partial x_k} = \frac{\partial}{\partial x_k} (x_i x_j) = \delta_{ik} x_j + \delta_{jk} x_i,
    \]
    где $\delta_{ik}$ — символ Кронекера:
    \[
    \delta_{ik} =
    \begin{cases}
    1, & \text{если } i = k, \\
    0, & \text{иначе}.
    \end{cases}
    \]

    Якобиан в покомпонентной форме:
    Строка якобиана, соответствующая элементу $(i, j)$, имеет вид:
    \[
    J_{f_3}^{(i,j)}(x) =
    \begin{bmatrix}
    \delta_{i1} x_j + \delta_{j1} x_i &
    \delta_{i2} x_j + \delta_{j2} x_i &
    \cdots &
    \delta_{in} x_j + \delta_{jn} x_i
    \end{bmatrix} \in \mathbb{R}^{1 \times n}.
    \]

    С использованием произведения Кронекера $\otimes$:
    \[
    \boxed{J_{f_3}(x) = \frac{\partial \mathrm{vec}(xx^\top)}{\partial x^\top} = I_n \otimes x^\top + x^\top \otimes I_n \in \mathbb{R}^{n^2 \times n}}.
    \]

    Пусть $I_n$ — единичная матрица размера $n \times n$, тогда:
    \[
    \begin{aligned}
    I_n \otimes x^\top &=
    \begin{bmatrix}
    x^\top & 0 & \cdots & 0 \\
    0 & x^\top & \cdots & 0 \\
    \vdots & \vdots & \ddots & \vdots \\
    0 & 0 & \cdots & x^\top
    \end{bmatrix} \in \mathbb{R}^{n^2 \times n}, \\
    x^\top \otimes I_n &=
    \begin{bmatrix}
    x_1 I_n \\
    x_2 I_n \\
    \vdots \\
    x_n I_n
    \end{bmatrix} \in \mathbb{R}^{n^2 \times n}.
    \end{aligned}
    \]

    \textbf{Ответ}:
    \begin{enumerate}[label=\alph*)]
        \item
        \[
        \boxed{\frac{\partial f_1}{\partial x} \in \mathbb{R}^{1 \times 2}}.
        \]
        \[
        \boxed{\frac{\partial f_2}{\partial x} \in \mathbb{R}^{1 \times n}}.
        \]
        \[
        \boxed{\frac{\partial f_3}{\partial x} \equiv \frac{\partial \mathrm{vec}(xx^\top)}{\partial x} \in \mathbb{R}^{n^2 \times n}}.
        \]
        \item
        \[
        \boxed{J_{f_1}(x) = \frac{\partial f_1}{\partial x} =
        \begin{bmatrix}
        \cos x_1 \cos x_2 & -\sin x_1 \sin x_2
        \end{bmatrix}}.
        \]
        \[
        \boxed{J_{f_2}(x, y) = \frac{\partial f_2}{\partial x} =
        \begin{bmatrix}
        y_1 & y_2 & \cdots & y_n
        \end{bmatrix} = y^\top}.
        \]
        \[
            \boxed{
        J_{f_3}^{(i,j)}(x) =
        \begin{bmatrix}
        \delta_{i1} x_j + \delta_{j1} x_i &
        \delta_{i2} x_j + \delta_{j2} x_i &
        \cdots &
        \delta_{in} x_j + \delta_{jn} x_i
        \end{bmatrix}
            }
        \]
    \end{enumerate}

    \textbf{Условие задания 5.6}:

    Вычислить производные следующих функций:
    \begin{enumerate}
        \item[$\bullet$] $f(t) = \sin(\log(t^\top t))$, где $t \in \mathbb{R}^D$
        \item[$\bullet$] $g(X) = \operatorname{tr}(AXB)$, где $A \in \mathbb{R}^{D \times E}$, $X \in \mathbb{R}^{E \times F}$, $B \in \mathbb{R}^{F \times D}$
    \end{enumerate}

    \textbf{Решение}:

    Решение для функции $f(t)$:
    Введём вспомогательную функцию:
    \[
    s(t) = t^\top t = \sum_{i=1}^D t_i^2.
    \]

    Тогда:
    \[
    f(t) = \sin(\log s(t)).
    \]

    Применяем цепное правило:

    \begin{enumerate}
        \item[$\bullet$] Производная внешней функции:
        \[
        \frac{d}{ds} \sin(\log s) = \cos(\log s) \cdot \frac{d}{ds} \log s
        = \cos(\log s) \cdot \frac{1}{s}.
        \]

        \item[$\bullet$] Градиент $s(t)$:
        \[
        \nabla_t s(t) = \frac{\partial}{\partial t} (t^\top t) = 2t.
        \]

        \item[$\bullet$] По цепному правилу:
        \[
        \nabla_t f(t) = \cos(\log s(t)) \cdot \frac{1}{s(t)} \cdot \nabla_t s(t)
        = \cos(\log(t^\top t)) \cdot \frac{1}{t^\top t} \cdot 2t.
        \]
    \end{enumerate}

    \[
    \boxed{\dfrac{\partial f}{\partial t} = \dfrac{2\cos\left(\log(t^\top t)\right)}{t^\top t} \cdot t \in \mathbb{R}^D}
    \]

    В матричной форме (якобиан как строка):
    \[
    \boxed{\dfrac{\partial f}{\partial t^\top} = \dfrac{2\cos\left(\log(t^\top t)\right)}{t^\top t} \cdot t^\top \in \mathbb{R}^{1 \times D}}
    \]

    Решение для функции $g(X)$:
    Используем свойство цикличности следа:
    \[
    \operatorname{tr}(AXB) = \operatorname{tr}(BAX).
    \]

    Обозначим $C = BA \in \mathbb{R}^{F \times E}$.

    Тогда:
    \[
    g(X) = \operatorname{tr}(CX).
    \]

    Известная формула для производной следа:
    \[
    \frac{\partial}{\partial X} \operatorname{tr}(CX) = C^\top.
    \]

    Подставляем $C = BA$:
    \[
    \frac{\partial g}{\partial X} = (BA)^\top = A^\top B^\top.
    \]

    Проверяем размерности:
    \begin{itemize}
        \item[$\bullet$] $A^\top \in \mathbb{R}^{E \times D}$
        \item[$\bullet$] $B^\top \in \mathbb{R}^{D \times F}$
        \item[$\bullet$] $A^\top B^\top \in \mathbb{R}^{E \times F}$
    \end{itemize}

    Размерность совпадает с размерностью $X \in \mathbb{R}^{E \times F}$.

    \[
    \boxed{\dfrac{\partial g}{\partial X} = A^\top B^\top \in \mathbb{R}^{E \times F}}
    \]

    \textbf{Ответ}:
    \begin{enumerate}
        \item Для $f(t) = \sin(\log(t^\top t))$:
        \[
        \boxed{\dfrac{\partial f}{\partial t} = \dfrac{2\cos\left(\log(t^\top t)\right)}{t^\top t} \cdot t}
        \]

        \item Для $g(X) = \operatorname{tr}(AXB)$:
        \[
        \boxed{\dfrac{\partial g}{\partial X} = A^\top B^\top}
        \]
    \end{enumerate}

    \textbf{Условие задания 5.7}:

    Вычислить производные следующих функций, используя цепное правило в матричной форме:
    \begin{enumerate}[label=\alph*)]
        \item $f(z) = \log(1 + z)$, где $z = x^\top x$, $x \in \mathbb{R}^D$
        \item $f(z) = \sin(z)$, где $z = Ax + b$, $x \in \mathbb{R}^D$, $A \in \mathbb{R}^{E \times D}$, $b \in \mathbb{R}^E$
    \end{enumerate}

    \textbf{Решение}:

    Для композиции функций $y = g(z)$, $z = h(x)$:
    \[
    \frac{\partial y}{\partial x} = \frac{\partial y}{\partial z} \cdot \frac{\partial z}{\partial x},
    \]
    где размерности согласованы: если $y \in \mathbb{R}^p$, $z \in \mathbb{R}^q$, $x \in \mathbb{R}^m$, то:
    \[
    \frac{\partial y}{\partial x} \in \mathbb{R}^{p \times m}, \quad
    \frac{\partial y}{\partial z} \in \mathbb{R}^{p \times q}, \quad
    \frac{\partial z}{\partial x} \in \mathbb{R}^{q \times m}.
    \]

    \begin{itemize}
        \item[$\bullet$] $x \in \mathbb{R}^D$ — вектор-столбец
        \item[$\bullet$] $z = x^\top x \in \mathbb{R}$ — скаляр
        \item[$\bullet$] $f(z) \in \mathbb{R}$ — скаляр
    \end{itemize}

    Искомый якобиан: $\dfrac{\partial f}{\partial x} \in \mathbb{R}^{1 \times D}$.

    Производная $f$ по $z$:
    \[
    \frac{df}{dz} = \frac{d}{dz} \log(1 + z) = \frac{1}{1 + z} \in \mathbb{R}^{1 \times 1}.
    \]

    Производная $z$ по $x$:
    \[
    z = x^\top x = \sum_{i=1}^D x_i^2.
    \]

    Частные производные:
    \[
    \frac{\partial z}{\partial x_i} = 2x_i, \quad i = 1, \dots, D.
    \]

    В векторной форме (якобиан как строка):
    \[
    \frac{\partial z}{\partial x} = \begin{bmatrix}
    2x_1 & 2x_2 & \cdots & 2x_D
    \end{bmatrix} = 2x^\top \in \mathbb{R}^{1 \times D}.
    \]

    Применение цепного правила:
    \[
    \frac{\partial f}{\partial x} = \frac{df}{dz} \cdot \frac{\partial z}{\partial x}
    = \frac{1}{1 + z} \cdot 2x^\top
    = \frac{2}{1 + x^\top x} \cdot x^\top.
    \]

    \[
    \boxed{\frac{\partial f}{\partial x} = \frac{2}{1 + x^\top x} \cdot x^\top \in \mathbb{R}^{1 \times D}}
    \]

    В форме градиента (вектор-столбец):
    \[
    \boxed{\nabla_x f = \frac{2}{1 + x^\top x} \cdot x \in \mathbb{R}^{D \times 1}}
    \]

    \begin{itemize}
        \item[$\bullet$] $x \in \mathbb{R}^D$ — вход
        \item[$\bullet$] $A \in \mathbb{R}^{E \times D}$ — матрица
        \item[$\bullet$] $b \in \mathbb{R}^E$ — вектор смещения
        \item[$\bullet$] $z = Ax + b \in \mathbb{R}^E$ — промежуточный вектор
        \item[$\bullet$] $f(z) \in \mathbb{R}^E$ — выход
    \end{itemize}

    Искомый якобиан: $\dfrac{\partial f}{\partial x} \in \mathbb{R}^{E \times D}$.

    Производная $f$ по $z$. Для каждого компонента:
    \[
    f_i(z) = \sin(z_i), \quad i = 1, \dots, E.
    \]

    Частные производные:
    \[
    \frac{\partial f_i}{\partial z_j} =
    \begin{cases}
    \cos(z_i), & \text{если } i = j, \\
    0, & \text{если } i \neq j.
    \end{cases}
    \]

    В матричной форме:
    \[
    \frac{\partial f}{\partial z} =
    \begin{bmatrix}
    \cos(z_1) & 0 & \cdots & 0 \\
    0 & \cos(z_2) & \cdots & 0 \\
    \vdots & \vdots & \ddots & \vdots \\
    0 & 0 & \cdots & \cos(z_E)
    \end{bmatrix}
    = \operatorname{diag}\big(\cos(z)\big) \in \mathbb{R}^{E \times E}.
    \]

    Производная $z$ по $x$:
    \[
    z = Ax + b.
    \]

    Компонентно:
    \[
    z_i = \sum_{j=1}^D A_{ij} x_j + b_i.
    \]

    Частные производные:
    \[
    \frac{\partial z_i}{\partial x_j} = A_{ij}.
    \]

    В матричной форме:
    \[
    \frac{\partial z}{\partial x} = A \in \mathbb{R}^{E \times D}.
    \]

    Применение цепного правила:
    \[
    \frac{\partial f}{\partial x} = \frac{\partial f}{\partial z} \cdot \frac{\partial z}{\partial x}
    = \operatorname{diag}\big(\cos(z)\big) \cdot A.
    \]

    Подставляя $z = Ax + b$:
    \[
    \frac{\partial f}{\partial x} = \operatorname{diag}\big(\cos(Ax + b)\big) \cdot A.
    \]

    Покомпонентно:
    \[
    \left(\frac{\partial f}{\partial x}\right)_{ij} = \cos\big((Ax + b)_i\big) \cdot A_{ij}.
    \]

    \[
    \boxed{\frac{\partial f}{\partial x} = \operatorname{diag}\big(\cos(Ax + b)\big) \cdot A \in \mathbb{R}^{E \times D}}
    \]

    \textbf{Ответ}:
    \item[$\bullet$]
    \[
    \boxed{\frac{\partial f}{\partial x} = \frac{2}{1 + x^\top x} \cdot x^\top \in \mathbb{R}^{1 \times D}}
    \]
    \item[$\bullet$]
    \[
    \boxed{\frac{\partial f}{\partial x} = \operatorname{diag}\big(\cos(Ax + b)\big) \cdot A \in \mathbb{R}^{E \times D}}
    \]

    \textbf{Условие задания 5.8}:

    Вычислить производные следующих функций, используя цепное правило и указывая размерности всех промежуточных производных:
    \begin{enumerate}[label=\alph*)]
        \item $f(z) = \exp\left(-\frac{1}{2}z\right)$, где $z = g(y) = y^\top S^{-1}y$, $y = h(x) = x - \mu$,
              $x, \mu \in \mathbb{R}^D$, $S \in \mathbb{R}^{D \times D}$ — симметричная положительно определённая матрица

        \item $f(x) = \operatorname{tr}(xx^\top + \sigma^2 I)$, $x \in \mathbb{R}^D$, $\sigma \in \mathbb{R}$

        \item $f = \tanh(z)$, где $z = Ax + b$, $x \in \mathbb{R}^N$, $A \in \mathbb{R}^{M \times N}$, $b \in \mathbb{R}^M$
    \end{enumerate}

    \textbf{Решение}:

    Для композиции $f(z)$, $z = g(y)$, $y = h(x)$:
    \[
    \frac{\partial f}{\partial x} = \frac{\partial f}{\partial z} \cdot \frac{\partial z}{\partial y} \cdot \frac{\partial y}{\partial x},
    \]
    где умножение понимается как матричное.

    Если $f \in \mathbb{R}^p$, $z \in \mathbb{R}^q$, $y \in \mathbb{R}^r$, $x \in \mathbb{R}^s$, то:
    \[
    \frac{\partial f}{\partial x} \in \mathbb{R}^{p \times s}, \quad
    \frac{\partial f}{\partial z} \in \mathbb{R}^{p \times q}, \quad
    \frac{\partial z}{\partial y} \in \mathbb{R}^{q \times r}, \quad
    \frac{\partial y}{\partial x} \in \mathbb{R}^{r \times s}.
    \]

    \begin{itemize}
        \item[$\bullet$] $x \in \mathbb{R}^D$
        \item[$\bullet$] $y = x - \mu \in \mathbb{R}^D$
        \item[$\bullet$] $z = y^\top S^{-1}y \in \mathbb{R}$ (скаляр)
        \item[$\bullet$] $f(z) \in \mathbb{R}$ (скаляр)
    \end{itemize}

    Искомый якобиан: $\dfrac{\partial f}{\partial x} \in \mathbb{R}^{1 \times D}$.

    Производная $f$ по $z$:
    \[
    \frac{df}{dz} = \frac{d}{dz} \exp\left(-\frac{1}{2}z\right) = -\frac{1}{2} \exp\left(-\frac{1}{2}z\right) = -\frac{1}{2} f(z).
    \]

    Размерность: $\dfrac{df}{dz} \in \mathbb{R}^{1 \times 1}$.

    Производная $z$ по $y$:
    \[
    z = y^\top S^{-1}y = y^\top A y, \quad \text{где } A = S^{-1}.
    \]

    Для квадратичной формы с симметричной матрицей $A$:
    \[
    \frac{\partial z}{\partial y} = 2y^\top A = 2y^\top S^{-1}.
    \]

    Размерность: $\dfrac{\partial z}{\partial y} \in \mathbb{R}^{1 \times D}$.

    Производная $y$ по $x$:
    \[
    y = x - \mu \quad \Rightarrow \quad \frac{\partial y}{\partial x} = I_D,
    \]
    где $I_D$ — единичная матрица размера $D \times D$.

    Размерность: $\dfrac{\partial y}{\partial x} \in \mathbb{R}^{D \times D}$.

    Применение цепного правила:
    \[
    \frac{\partial f}{\partial x} = \frac{df}{dz} \cdot \frac{\partial z}{\partial y} \cdot \frac{\partial y}{\partial x}
    = \left(-\frac{1}{2} f\right) \cdot \left(2y^\top S^{-1}\right) \cdot I_D
    = -f \cdot y^\top S^{-1}.
    \]

    Подставляем выражения для $f$ и $y$:
    \[
    f(z) = \exp\left(-\frac{1}{2} (x - \mu)^\top S^{-1} (x - \mu)\right), \quad y = x - \mu.
    \]

    \[
    \boxed{\frac{\partial f}{\partial x} = -\exp\left(-\frac{1}{2} (x - \mu)^\top S^{-1} (x - \mu)\right) \cdot (x - \mu)^\top S^{-1} \in \mathbb{R}^{1 \times D}}
    \]

    В форме градиента (вектор-столбец):
    \[
    \boxed{\nabla_x f = -\exp\left(-\frac{1}{2} (x - \mu)^\top S^{-1} (x - \mu)\right) \cdot S^{-1} (x - \mu) \in \mathbb{R}^{D \times 1}}
    \]

    Используем свойства следа:
    \[
    \operatorname{tr}(xx^\top) = x^\top x, \quad \operatorname{tr}(\sigma^2 I) = \sigma^2 \operatorname{tr}(I) = \sigma^2 D.
    \]

    Тогда:
    \[
    f(x) = x^\top x + \sigma^2 D.
    \]

    Константа $\sigma^2 D$ не зависит от $x$, поэтому:
    \[
    \frac{\partial f}{\partial x} = \frac{\partial}{\partial x} (x^\top x).
    \]

    Для квадратичной формы:
    \[
    \frac{\partial}{\partial x} (x^\top x) = 2x^\top.
    \]

    \[
    \boxed{\frac{\partial f}{\partial x} = 2x^\top \in \mathbb{R}^{1 \times D}}
    \]

    В форме градиента:
    \[
    \boxed{\nabla_x f = 2x \in \mathbb{R}^{D \times 1}}
    \]

    \begin{itemize}
        \item[$\bullet$] $x \in \mathbb{R}^N$
        \item[$\bullet$] $z = Ax + b \in \mathbb{R}^M$
        \item[$\bullet$] $f = \tanh(z) \in \mathbb{R}^M$
    \end{itemize}

    Искомый якобиан: $\dfrac{\partial f}{\partial x} \in \mathbb{R}^{M \times N}$.

    Производная $f$ по $z$. Для каждого компонента:
    \[
    f_i(z) = \tanh(z_i), \quad i = 1, \dots, M.
    \]

    Из одномерного анализа:
    \[
    \frac{d}{dz} \tanh(z) = 1 - \tanh^2(z).
    \]

    Поэтому:
    \[
    \frac{\partial f_i}{\partial z_j} =
    \begin{cases}
    1 - \tanh^2(z_i), & \text{если } i = j, \\
    0, & \text{если } i \neq j.
    \end{cases}
    \]

    В матричной форме:
    \[
    \frac{\partial f}{\partial z} = \operatorname{diag}\bigl(1 - \tanh^2(z)\bigr) \in \mathbb{R}^{M \times M}.
    \]

    Производная $z$ по $x$:
    \[
    z = Ax + b \quad \Rightarrow \quad \frac{\partial z}{\partial x} = A \in \mathbb{R}^{M \times N}.
    \]

    Применение цепного правила:
    \[
    \frac{\partial f}{\partial x} = \frac{\partial f}{\partial z} \cdot \frac{\partial z}{\partial x}
    = \operatorname{diag}\bigl(1 - \tanh^2(z)\bigr) \cdot A.
    \]

    Подставляя $z = Ax + b$:
    \[
    \frac{\partial f}{\partial x} = \operatorname{diag}\bigl(1 - \tanh^2(Ax + b)\bigr) \cdot A.
    \]

    \[
    \boxed{\frac{\partial f}{\partial x} = \operatorname{diag}\bigl(1 - \tanh^2(Ax + b)\bigr) \cdot A \in \mathbb{R}^{M \times N}}
    \]

    \textbf{Ответ}:
    \begin{enumerate}[label=\alph*)]
        \item
        \[
        \boxed{\frac{\partial f}{\partial x} = \frac{2}{1 + x^\top x} \cdot x^\top \in \mathbb{R}^{1 \times D}}
        \]
        \item
        \[
            \boxed{
        \frac{\partial}{\partial x} (x^\top x) = 2x^\top
            }
        \]
        \item
        \[
        \boxed{\frac{\partial f}{\partial x} = \operatorname{diag}\bigl(1 - \tanh^2(Ax + b)\bigr) \cdot A \in \mathbb{R}^{M \times N}}
        \]
    \end{enumerate}

    \textbf{Условие задания 5.9}:

    Дана функция:
    \[
    g(x, z, \nu) = \log p(x, z) - \log q(z, \nu),
    \]
    где $z$ зависит от $\nu$ через параметрическое преобразование:
    \[
    z = t(\varepsilon, \nu).
    \]

    Параметры:
    \begin{itemize}
        \item[$\bullet$] $x \in \mathbb{R}^D$ — фиксированный вектор (данные)
        \item[$\bullet$] $z \in \mathbb{R}^E$ — скрытая переменная
        \item[$\bullet$] $\nu \in \mathbb{R}^F$ — параметры вариационного распределения
        \item[$\bullet$] $\varepsilon \in \mathbb{R}^G$ — вспомогательная случайная переменная
        \item[$\bullet$] Функции $p(x, z)$, $q(z, \nu)$, $t(\varepsilon, \nu)$ — дифференцируемы по всем аргументам
    \end{itemize}

    Требуется вычислить градиент по $\nu$:
    \[
    \nabla_\nu g = \nabla_\nu \bigl[ \log p(x, t(\varepsilon, \nu)) - \log q(t(\varepsilon, \nu), \nu) \bigr].
    \]

    \textbf{Решение}:

    Рассмотрим первое слагаемое:
    \[
    \phi(\nu) = \log p(x, t(\varepsilon, \nu)).
    \]

    Функция $\phi$ зависит от $\nu$ только через $z = t(\varepsilon, \nu)$. По правилу цепочки:

    \[
    \nabla_\nu \phi(\nu) = \left( \frac{\partial t}{\partial \nu} \right)^\top \cdot \nabla_z \log p(x, z) \big|_{z = t(\varepsilon, \nu)}.
    \]

    \begin{itemize}
        \item[$\bullet$] $\nabla_z \log p(x, z) \in \mathbb{R}^E$ — градиент по $z$
        \item[$\bullet$] $\dfrac{\partial t}{\partial \nu} \in \mathbb{R}^{E \times F}$ — якобиан
        \item[$\bullet$] $\left(\dfrac{\partial t}{\partial \nu}\right)^\top \cdot \nabla_z \log p(x, z) \in \mathbb{R}^F$
    \end{itemize}

    Рассмотрим второе слагаемое:
    \[
    \psi(\nu) = \log q(t(\varepsilon, \nu), \nu).
    \]

    Функция $\psi$ зависит от $\nu$ двумя способами:
    \begin{enumerate}
        \item[$\bullet$] Через $z = t(\varepsilon, \nu)$ (косвенная зависимость)
        \item[$\bullet$] Непосредственно через второй аргумент (прямая зависимость)
    \end{enumerate}

    Применяем правило цепочки для нескольких путей:

    \[
    \nabla_\nu \psi(\nu) = \underbrace{\left( \frac{\partial t}{\partial \nu} \right)^\top \cdot \nabla_z \log q(z, \nu)}_{\text{зависимость через } z} + \underbrace{\nabla_\nu \log q(z, \nu)}_{\text{прямая зависимость}} \big|_{z = t(\varepsilon, \nu)}.
    \]

    \begin{itemize}
        \item[$\bullet$] Первое слагаемое: $\left(\dfrac{\partial t}{\partial \nu}\right)^\top \cdot \nabla_z \log q(z, \nu) \in \mathbb{R}^F$
        \item[$\bullet$] Второе слагаемое: $\nabla_\nu \log q(z, \nu) \in \mathbb{R}^F$
        \item[$\bullet$] Сумма: $\nabla_\nu \psi(\nu) \in \mathbb{R}^F$
    \end{itemize}

    Градиент полной функции:
    \[
    \nabla_\nu g = \nabla_\nu \phi(\nu) - \nabla_\nu \psi(\nu).
    \]

    Подставляем вычисленные выражения:
    \[
    \begin{aligned}
    \nabla_\nu g(x, z, \nu) &= \left( \frac{\partial t}{\partial \nu} \right)^\top \cdot \nabla_z \log p(x, z) \\
    &\quad - \left[ \left( \frac{\partial t}{\partial \nu} \right)^\top \cdot \nabla_z \log q(z, \nu) + \nabla_\nu \log q(z, \nu) \right] \\
    &= \left( \frac{\partial t}{\partial \nu} \right)^\top \cdot \left[ \nabla_z \log p(x, z) - \nabla_z \log q(z, \nu) \right] - \nabla_\nu \log q(z, \nu),
    \end{aligned}
    \]

    где $z = t(\varepsilon, \nu)$.

    \textbf{Ответ}:
    \[
    \boxed{
    \begin{aligned}
    \nabla_\nu g(x, z, \nu) &= \left( \frac{\partial t}{\partial \nu} \right)^\top \cdot \left[ \nabla_z \log p(x, z) - \nabla_z \log q(z, \nu) \right] \\
    &\quad - \nabla_\nu \log q(z, \nu), \quad \text{где } z = t(\varepsilon, \nu).
    \end{aligned}
    }
    \]























    \newpage



    \appendix

\end{document}

